\uuid{2Jhe}
\titre{Raisonnements}
\auteur{bodin ; croizat ; sacre}
\organisation{exo7}
\licence{CC BY-NC-SA 4.0}
\datecreate{2022-03-08}
\niveau{L1}
\module{Algèbre}
\chapitre{Logique, ensemble, raisonnement}
\sousChapitre{Absurde et contraposée}
\difficulte{2}
\theme{Raisonnements}
\exo7qcmid{hilisit.logique.047}

\contenu{
\begin{question}
\qclassification{100.03}
Quel type de raisonnement est adapté pour montrer qu'il existe une infinité de nombres premiers ?
\begin{answers}
	\bad{Au cas par cas : on étudie $n=2$, $n=3$, $n=5$,...}
	\bad{Par récurrence sur $n$ parcourant l'ensemble des nombres premiers.}
	\good{Par l'absurde en supposant qu'il n'existe qu'un nombre fini de nombres premiers.}
	\bad{C'est une propriété que l'on ne sait pas démontrer.}
\end{answers}
\begin{explanations}
La démonstration classique se fait par l'absurde en supposant qu'il n'existe qu'un nombre fini de nombres premiers puis en cherchant une contradiction.
\end{explanations}
\end{question}
}
